%!TEX program = pdflatex
\documentclass[UTF8]{ctexart}

\usepackage[a4paper, scale=.8]{geometry}
\usepackage{listings}
\usepackage{amsmath}


\pagestyle{empty}
\ctexset{
    section/format+ = \raggedright
}

\title{论文思路}
\author{xxx}

\begin{document}
\maketitle
\section{关键词}
% 描述研究领域、方法、技术等的专业词汇

\section{背景}
\subsection{应用场景}
% 使用文字或图像描绘论文的应用场景，建议150字以下

\subsection{问题陈述}
% 陈述在该场景中存在的安全、隐私、效率等问题，建议150字以下
安全协议的形式化分析工作中缺少通用的自动化的密钥暴力破解的分析方法。

\subsection{相关工作}
% 分析与论文场景或问题相关的、近三年或经典的研究的不足，并附上相应的参考文献

\section{意义}
\subsection{主要挑战}
% 阐述论文方案面临的挑战，建议200字以下

\begin{enumerate}
    \item 建模方式通用化，所有存在暴力破解问题的协议都能够使用该方法进行建模。
    \item 拓展 Tamarin-Prover 语法，密钥类型化，优化模型验证流程。
    \item 分析密钥派生过程，自动化生成密钥信息熵继承、迁移规则。
    \item 密钥暴力破解建模不影响原有模型的验证功能和验证速度
\end{enumerate}


\subsection{论文贡献}
% 阐述论文的主要贡献，建议200字以下

\begin{enumerate}
    \item 通用的自动化的安全协议密钥建模方法，发现潜在的可被爆破的脆弱密钥。
    \item 易于使用的类型化的 Tamarin-Prover 语法拓展，模块化密钥建模。
    \item 优化协议模型结构，加快验证速度。
\end{enumerate}

\section{方案}
\subsection{安全模型}
% 描述敌手能力和目的

\subsection{技术依据}
% 介绍论文使用的关键技术及其核心思想，并附上相应的参考文献

\subsection{方案描述}
% 陈述论文的具体方案

\subsubsection{密钥熵迁移规则}

从攻击者的角度，爆破密钥的难度取决于该密钥所有的随机因素的遍历空间规模。

信息熵的可加性：如果子系统之间的相互作用是已知的，可以通过子系统的熵来计算一个系统的熵。

根据信息熵的可加性，设计四种密钥的熵值：
\begin{itemize}
    \item 零熵（Zero Entropy）：攻击者已知信息的熵。
    \item 低熵（Low Entropy）：攻击者可以通过暴力破解获取的最长密钥的熵值。\\
            $1\;Zero < 1\;Low < 1\;High$
    \item 高熵（High Entropy）：攻击者直接通过暴力破解的方式难以获取的密钥的熵值。\\
            $1\;Low < 1\;High < 1\;Chaos$
    \item 混沌熵（Chaos Entropy）：攻击者永远无法通过暴力破解获取的密钥的熵值。\\
            $1\;High < 1\;Chaos$
\end{itemize}

根据信息熵的可加性列举出以下规则，用于描述信息熵在密钥派生过程中的传递。

熵增：
\begin{itemize}
    \item Zero + X = X
    \item Low + Low = Low (completeness)
    \item Low + Low = High (soundness)
    \item Low + High = High
    \item High + High = High
    \item Chaos + X = Chaos
\end{itemize}

熵减：
\begin{itemize}
    \item $<X, Y> - X = Y$
\end{itemize}

\subsubsection{文法描述}
% \begin{lstlisting}
%     rule        :=  simple_rule [variants]
%     diff_rule   :=  simple_rule ['left' rule 'right' rule]
%     simple_rule :=  'rule' [modulo] ident [rule_attrs] [key_block]':'
%                     [let_block]
%                     '[' [facts] ']' ( '-->' | '--[' facts ']->') '[' [facts] ']'
%     variants    :=  'variants' simple_rule (',' simple_rule)*
%     modulo      :=  '(' 'modulo' ('E' | 'AC') ')'
%     rule_attrs  :=  '[' rule_attr (',' rule_attr)* [','] ']'
%     rule_attr   :=  ('color=' | 'colour=') hexcolor
%     let_block   :=  'let' (msg_var '=' msetterm)+ 'in'
%     msg_var     :=  ident ['.' natural] [':' 'msg']
%     hexcolor    :=  "'" ["#"] hexdigit{6} "'" | ["#"] hexdigit{6}
%     ident       :=  alphaNum (alphaNum | '_')*

%     key_block   :=  '('  ')'
%     key_type    :=  'low' | 'medium' | 'high' | 'chaotic'

%     facts := fact (',' fact)*
%     fact  := ['!'] ident '(' [msetterm (',' msetterm)*] ')' [fact_annotes]
%     fact_annotes := '[' fact_annote (',' fact_annote)* ']'
%     fact_annote  := '+' | '-' | 'no_precomp'

%     literal     := "'"  ~('\'' | '\n')+ "'" // a fixed, public name
%              | "~'" ~('\'' | '\n')+ "'" // a fixed, fresh name
%              | nonnode_var    // a non-temporal variable
%     nonnode_var := ['$'] ident ['.' natural]         // 'pub' sort prefix
%                 | ident ['.' natural] ':' 'pub'     // 'pub' sort suffix
%                 | ['~'] ident ['.' natural]         // 'fresh' sort prefix
%                 | ident ['.' natural] ':' 'fresh'   // 'fresh' sort suffix
%                 | msg_var                                // 'msg' sort
% \end{lstlisting}

% \begin{flalign*}
%     ident\; :=\; &alphaNum\; (alphaNum\; |\; '\_')^*\\
%     rule\;  :=\; &simple\_rule [variants]\\
%     simple\_rule\; :=\;  &\text{`rule'}\; [modulo]\; ident\; [rule\_attrs]\; \text{`:'}\\
%                 &[let\_block]\\
%                 &'['\; [facts]\; ']'\; (\; '-->'\; |\; '--[' facts ']->')\; '['\; [facts] ']'\\
% \end{flalign*}


Example

\begin{lstlisting}
    rule Register_pk(~lkt := chaos 'ltk'):
    [ Fr(~ltk) ]
  -->
    [ !Ltk($A, ~ltk), !Pk($A, pk(~ltk)) ] 

    rule Reveal_ltk(ltk = 'ltk'):
    [ !Ltk(A, ltk) ]
  --[ LtkReveal(A) ]->
    [ Out(ltk) ]

\end{lstlisting}


\section{评估}
\subsection{评价标准}
% 列出评价论文方案的理论或实验标准

\begin{itemize}
    \item 通用性：
        \begin{itemize}
            \item BC: 密钥降级
            \item BLE：密钥降级、长度混淆
            \item WPA 口令离线字典攻击
        \end{itemize}
    \item 效率：对比原模型文件修改前后的验证时间
\end{itemize}

    

\subsection{对比方案}
% 列出拟对比的方案（若为空则删除本小节），并附上相应的参考文献

\subsection{实验设置}
% 描述包含哪些实验，每个实验的数据集和步骤（若无实验则删除本小节）

\bibliographystyle{plain}
\bibliography{library}

\end{document}